\documentclass[10pt,a4paper]{amsart}
\usepackage{amsmath,amssymb,amsthm,graphicx}
% Retained minimal preamble for the Stacks Project chapter more-algebra.tex (More on Algebra).
% The chapter file pinned-original-source.tex carries only \input{preamble} on its first line, so no
% preamble text lives inside the chapter itself.  The Stacks Project uses one single shared preamble for
% every chapter; the snapshot stages it identically as sources/stacks-chow/preamble.tex and
% sources/stacks-etale-cohomology/preamble.tex.  The block below is copied mechanically, line for line,
% from sources/stacks-chow/preamble.tex lines 14-18 (the Xy-pic setup) and 154-217 (the native theorem
% declarations together with the mathematical macros).  Nothing was added, removed or redefined; the
% wrapper class and the amsmath/amssymb/amsthm/graphicx base packages are supplied by the assembler.
% Chapter: more-algebra.tex -- More on Algebra.
% For commutative diagrams we use Xy-pic
\usepackage[all]{xy}

% We use 2cell for 2-commutative diagrams.
\xyoption{2cell}

% Theorem environments.
%
\theoremstyle{plain}
\newtheorem{theorem}[subsection]{Theorem}
\newtheorem{proposition}[subsection]{Proposition}
\newtheorem{lemma}[subsection]{Lemma}

\theoremstyle{definition}
\newtheorem{definition}[subsection]{Definition}
\newtheorem{example}[subsection]{Example}
\newtheorem{exercise}[subsection]{Exercise}
\newtheorem{situation}[subsection]{Situation}

\theoremstyle{remark}
\newtheorem{remark}[subsection]{Remark}
\newtheorem{remarks}[subsection]{Remarks}

\numberwithin{equation}{subsection}

% Macros
%
\def\lim{\mathop{\mathrm{lim}}\nolimits}
\def\colim{\mathop{\mathrm{colim}}\nolimits}
\def\Spec{\mathop{\mathrm{Spec}}}
\def\Hom{\mathop{\mathrm{Hom}}\nolimits}
\def\Ext{\mathop{\mathrm{Ext}}\nolimits}
\def\SheafHom{\mathop{\mathcal{H}\!\mathit{om}}\nolimits}
\def\SheafExt{\mathop{\mathcal{E}\!\mathit{xt}}\nolimits}
\def\Sch{\mathit{Sch}}
\def\Mor{\mathop{\mathrm{Mor}}\nolimits}
\def\Ob{\mathop{\mathrm{Ob}}\nolimits}
\def\Sh{\mathop{\mathit{Sh}}\nolimits}
\def\NL{\mathop{N\!L}\nolimits}
\def\CH{\mathop{\mathrm{CH}}\nolimits}
\def\proetale{{pro\text{-}\acute{e}tale}}
\def\etale{{\acute{e}tale}}
\def\QCoh{\mathit{QCoh}}
\def\Ker{\mathop{\mathrm{Ker}}}
\def\Im{\mathop{\mathrm{Im}}}
\def\Coker{\mathop{\mathrm{Coker}}}
\def\Coim{\mathop{\mathrm{Coim}}}

% Boxtimes
%
\DeclareMathSymbol{\boxtimes}{\mathbin}{AMSa}{"02}

%
% Macros for moduli stacks/spaces
%
\def\QCohstack{\mathcal{QC}\!\mathit{oh}}
\def\Cohstack{\mathcal{C}\!\mathit{oh}}
\def\Spacesstack{\mathcal{S}\!\mathit{paces}}
\def\Quotfunctor{\mathrm{Quot}}
\def\Hilbfunctor{\mathrm{Hilb}}
\def\Curvesstack{\mathcal{C}\!\mathit{urves}}
\def\Polarizedstack{\mathcal{P}\!\mathit{olarized}}
\def\Complexesstack{\mathcal{C}\!\mathit{omplexes}}
% \Pic is the operator that assigns to X its picard group, usage \Pic(X)
% \Picardstack_{X/B} denotes the Picard stack of X over B
% \Picardfunctor_{X/B} denotes the Picard functor of X over B
\def\Pic{\mathop{\mathrm{Pic}}\nolimits}
\def\Picardstack{\mathcal{P}\!\mathit{ic}}
\def\Picardfunctor{\mathrm{Pic}}
\def\Deformationcategory{\mathcal{D}\!\mathit{ef}}
\usepackage{cleveref}
\crefname{originaltheorem}{original theorem}{original theorems}
\Crefname{originaltheorem}{Original theorem}{Original theorems}
\makeatletter
\newlabel{algebra-theorem-projective-free-over-local-ring}{{0593}{source}{}{sourcepointer-1}{}}
\newlabel{algebra-theorem-projective-free-over-local-ring@cref}{{[originaltheorem][0][]0593}{[1][0][]source}{}{}{}}
\makeatother
\crefname{originallemma}{original lemma}{original lemmas}
\Crefname{originallemma}{Original lemma}{Original lemmas}
\makeatletter
\newlabel{algebra-lemma-contained-in-radical}{{0AME}{source}{}{sourcepointer-2}{}}
\newlabel{algebra-lemma-contained-in-radical@cref}{{[originallemma][0][]0AME}{[1][0][]source}{}{}{}}
\makeatother
\crefname{originallemma}{original lemma}{original lemmas}
\Crefname{originallemma}{Original lemma}{Original lemmas}
\makeatletter
\newlabel{algebra-lemma-tor-noetherian}{{0AZ4}{source}{}{sourcepointer-3}{}}
\newlabel{algebra-lemma-tor-noetherian@cref}{{[originallemma][0][]0AZ4}{[1][0][]source}{}{}{}}
\makeatother
\title{Correcting an element by the radical quotient to make a single generated submodule a free direct summand}
\author{The Stacks Project Authors}
\date{stacks-project:more-algebra.tex:local-snapshot-20261012; GNU FDL 1.2 or later}
\begin{document}
\maketitle
\setcounter{subsection}{1}
\noindent\textit{Editorial note.} All delivered mathematical text is the author's own, copied from the pinned Stacks Project chapter file more-algebra.tex; no statement or proof step was rewritten, reordered or removed. The counted statement is the labeled lemma environment at source lines 39025-39033 (canonical tag 0GVI) and the counted proof is the proof environment at source lines 39035-39105, which immediately follows it and is its own proof.\par\medskip
\noindent\textit{Editorial note.} The proof displays a defect in the author's text at proof lines 39042-39043, where the free direct summand is asserted of M' rather than of the fibre of P; the printed text is copied exactly as it stands and is not corrected, since the lane delivers the author's own words.\par\medskip
\noindent\textit{Editorial note.} The three references are all to labelled results of another chapter: freeness of projectives over local rings, the unit criterion modulo the radical, and finiteness of a Tor group for finite modules over a Noetherian ring. All three are navigation stubs with their roles recorded in the reasons.\par\medskip
\noindent\textit{Editorial note.} The chapter file pinned-original-source.tex carries only an input of the shared preamble on its first line, so no preamble text lives inside the chapter. The retained file sources/stacks-more-algebra/retained-preamble.tex is a mechanical, line-for-line copy of the project's single shared preamble as staged in the snapshot (Xy-pic setup plus the native theorem declarations and mathematical macros). Nothing was added, removed or redefined, and the excerpt serves only as preamble, never as a span coordinate.\par\medskip
\noindent\textit{Editorial note.} The source URL is the source-level Stacks Project URL, as required when more than one card is produced from one source; the per-claim tag is carried by claim\_key and by this note. The counted statement and proof cite no bibliography entry, so this card needs no bibliography.\par\medskip
\noindent\textit{Allowed standard prerequisites (editorial).}\begin{itemize}
\item A finitely generated projective module over a local ring is free, so the localizations of P at maximal ideals are free and the statement makes sense. invoked at proof lines 39036-39037. Makes the rank and freeness language meaningful for the localizations appearing in the hypotheses.
\item An element of a ring which is congruent to one modulo the Jacobson radical is a unit. invoked at proof lines 39049-39051. Turns the congruence of the dual functional modulo the radical into a splitting over the whole ring.
\item For a finite module over a Noetherian ring the first Tor group against a residue field is finite, and a finite dimensional vector space cannot surject onto a subquotient that has infinite dimension. used at proof lines 39069-39076. Produces an element t independent of the image of s + m in the fibre, which is the input to the next linear algebra step.
\end{itemize}\medskip
\ifdefined\hypertarget\hypertarget{sourcepointer-1}{}\fi
\noindent\textit{Source reference.} Original theorem 0593: sources/stacks-chow/tags.txt, source lines 6801--6801. The reference is invoked at proof lines 39036-39037 so that the localizations of P at maximal ideals are free and the statement makes sense; the construction that follows is inside the counted span.\par\medskip
\ifdefined\hypertarget\hypertarget{sourcepointer-2}{}\fi
\noindent\textit{Source reference.} Original lemma 0AME: sources/stacks-chow/tags.txt, source lines 13402--13402. The reference supplies the unit criterion modulo the Jacobson radical, used at proof lines 39049-39051 to turn the congruence of the lifted functional into a splitting; the lifting and the choice of m are inside the counted span.\par\medskip
\ifdefined\hypertarget\hypertarget{sourcepointer-3}{}\fi
\noindent\textit{Source reference.} Original lemma 0AZ4: sources/stacks-chow/tags.txt, source lines 13812--13812. The reference makes the kernel of the displayed map a quotient of a finite Tor group, used at proof lines 39072-39074 to conclude that the image of K in the fibre has infinite dimension; the linear algebra step built on that claim is inside the counted span.\par\medskip
\begin{lemma}
\label{lemma-trick-to-find-good-element}
Let $R$ be a ring with Jacobson radical $J$ such that $R/J$
is Noetherian. Let $P$ be a projective $R$-module
such that $P_\mathfrak m$ has infinite rank for all maximal ideals
$\mathfrak m$ of $R$. Let $s \in P$ and $M \subset P$ such that
$Rs + M = P$. Then we can find $m \in M$ such that $R(s + m)$
is a free direct summand of $P$.
\end{lemma}

\begin{proof}
The statement makes sense as $P_\mathfrak m$ is free by
Algebra, Theorem \ref{algebra-theorem-projective-free-over-local-ring}.

\medskip\noindent
Denote $M' \subset P/JP$ the image of $M$ and $s' \in P/JP$ the
image of $s$. Observe that $R/J s' + M' = P/JP$.
Suppose we can find $m' \in M'$ such that
$R/J(s' + m')$ is a free direct summand of $M'$.
Choose $\varphi' : P/JP \to R/J$ which gives a splitting, i.e.,
we have $\varphi'(s' + m') = 1$ in $R/J$.
Then since $P$ is a projective $R$-module we can find
a lift $\varphi : P \to R$ of $\varphi'$.
Choose $m \in M$ mapping to $m'$.
Then $\varphi(s + m) \in R$ is congruent to $1$ modulo $J$ and hence
a unit in $R$ (Algebra, Lemma \ref{algebra-lemma-contained-in-radical}).
Whence $R(s + m)$ is a free direct summand of $P$.
This reduces us to the case discussed in the next paragraph.

\medskip\noindent
Assume $R$ is Noetherian. Let $m \in M$ be an element and let
$\varphi_1, \ldots, \varphi_n : P \to R$ be $R$-linear maps. Denote
$$
Z(s + m, \varphi_1, \ldots, \varphi_n) \subset \Spec(R)
$$
the vanishing locus of $\varphi_1(s + m), \ldots, \varphi_n(s + m) \in R$.

\medskip\noindent
Suppose $\mathfrak m$ is a maximal ideal of $R$ and
$\mathfrak m \in Z(s + m, \varphi_1, \ldots, \varphi_n)$.
Set $K = M \cap \bigcap \Ker(\varphi_i)$. We claim the image of
$$
K/\mathfrak mK \to P/\mathfrak m P
$$
has infinite dimension. Namely, the quotient $P/K$ is a finite $R$-module
as it is isomorphic to a submodule of $P/M \oplus R^{\oplus n}$.
Thus we see that
the kernel of the displayed arrow is a quotient of
$\text{Tor}_1^R(P/K, \kappa(\mathfrak m))$ which is finite
by Algebra, Lemma \ref{algebra-lemma-tor-noetherian}.
Combined with the fact that $P/\mathfrak mP$ has infinite
dimension we obtain our claim. Thus we can find a $t \in K$
which maps to an element $\overline{t}$
of the vector space $P/\mathfrak mP$ which is linearly independent
of the image $\overline{s + m}$ of $s + m$.
By linear algebra, we find an $R$-linear map
$\overline{\varphi} : P \to \kappa(\mathfrak m)$
such that $\overline{\varphi}(\overline{t}) = 1$ and
$\overline{\varphi}(\overline{s + m}) = 0$.
Since $P$ is projective, we can find an $R$-linear map
$\varphi : P \to R$ lifting $\overline{\varphi}$.
Then we see that the vanishing locus
$Z(s + m + t, \varphi_1, \ldots, \varphi_n, \varphi)$ is
contained in $Z(s + m, \varphi_1, \ldots, \varphi_n)$ but
does not contain $\mathfrak m$, i.e., it
is strictly smaller than $Z(s + m, \varphi_1, \ldots, \varphi_n)$.

\medskip\noindent
Since $\Spec(R)$ is a Noetherian topological space, we see
from the arguments above that we may find
$m \in M$ and $\varphi_1, \ldots, \varphi_n : P \to R$
such that the closed subset $Z(s + m, \varphi_1, \ldots, \varphi_n)$
does not contain any closed points of $\Spec(R)$.
Hence $Z(s + m, \varphi_1, \ldots, \varphi_n) = \emptyset$.
Hence we can find $r_1, \ldots, r_n \in R$ such that
$\sum r_i\varphi_i(s + m) = 1$. Hence
$$
R \xrightarrow{s + m} P \xrightarrow{\sum r_i \varphi_i} R
$$
is the desired splitting.
\end{proof}

\end{document}
